\documentclass[12pt]{article}
\usepackage{amsmath, amssymb}
\usepackage{geometry}
\usepackage{enumitem}
\geometry{margin=1in}

\title{Polar Coordinates and Complex Polar Form Practice (Algebra II)}
\author{}
\date{}

\begin{document}
\maketitle

\section*{Instructions}
Complete in 60 minutes. Show clear work. Use exact values unless a decimal approximation is explicitly requested.

\section*{Relevant Formulas}
\begin{itemize}[leftmargin=1.5em]
  \item Polar to Cartesian: $x=r\cos\theta$, $y=r\sin\theta$
  \item Cartesian to Polar: $r^2=x^2+y^2$, $\tan\theta=\dfrac{y}{x}$
  \item Useful substitutions: $x^2+y^2=r^2$, $\dfrac{y}{x}=\tan\theta$, $r=\sqrt{x^2+y^2}$
  \item Polar form of a complex number: $z=r(\cos\theta+i\sin\theta)=r\,\text{cis}\,\theta$
  \item De Moivre's Formula: $\big[r(\cos\theta+i\sin\theta)\big]^n=r^n(\cos n\theta+i\sin n\theta)$
  \item Symmetry tests for polar curves $r=f(\theta)$:
  \begin{itemize}[leftmargin=1.5em]
    \item Polar axis: replace $\theta$ by $-\theta$
    \item Line $\theta=\dfrac{\pi}{2}$: replace $\theta$ by $\pi-\theta$
    \item Pole: replace $r$ by $-r$ or replace $\theta$ by $\theta+\pi$
  \end{itemize}
  \item Cartesian connection:
  \begin{itemize}[leftmargin=1.5em]
    \item Polar axis symmetry corresponds to $x$-axis symmetry
    \item Line $\theta=\dfrac{\pi}{2}$ symmetry corresponds to $y$-axis symmetry
    \item Pole symmetry corresponds to origin symmetry
  \end{itemize}
\end{itemize}

\section*{Problems}
\begin{enumerate}[leftmargin=*, itemsep=1em, topsep=0.8em]
  \item Convert the polar point $\left(-8,\frac{11\pi}{6}\right)$ to Cartesian coordinates.

  \item Convert the Cartesian point $(-4\sqrt{3},-4)$ to polar coordinates with $r>0$ and $0\le \theta<2\pi$.

  \item Find three additional polar representations of the point $\left(6,\frac{5\pi}{4}\right)$, including at least one with negative $r$.

  \item Convert the Cartesian equation
  \[
  x^2+y^2=6y
  \]
  to a polar equation.

  \item Convert the polar equation
  \[
  r=6\cos\theta-2\sin\theta
  \]
  to a Cartesian equation and identify the curve.

  \item Convert the Cartesian equation
  \[
  x^2+y^2=4x^2y
  \]
  to a polar equation.

  \item Convert the polar equation
  \[
  r^2=9\csc(2\theta)
  \]
  to a Cartesian equation.

  \item Test the polar curve
  \[
  r=2-5\sin\theta
  \]
  for symmetry about the polar axis, the line $\theta=\frac{\pi}{2}$, and the pole.

  \item Test the polar curve
  \[
  r^2=9\cos(2\theta)
  \]
  for symmetry about the polar axis, the line $\theta=\frac{\pi}{2}$, and the pole.

  \item Write the complex number $z=-3-3\sqrt{3}\,i$ in polar form $r(\cos\theta+i\sin\theta)$ with $0\le \theta<2\pi$.

  \item Use De Moivre's Formula to compute
  \[
  \left[3\left(\cos\frac{7\pi}{12}+i\sin\frac{7\pi}{12}\right)\right]^5
  \]
  in standard form.

  \item Use De Moivre's Formula to find all fourth roots of
  \[
  81\left(\cos\frac{5\pi}{6}+i\sin\frac{5\pi}{6}\right).
  \]
  Give each answer in polar form.

  \item Convert the Cartesian equation
  \[
  (x^2+y^2)^2=16xy
  \]
  to a polar equation, then state all three symmetry types the curve has.

  \item If
  \[
  z=2\left(\cos\frac{\pi}{9}+i\sin\frac{\pi}{9}\right),
  \]
  use De Moivre's Formula to find $z^6+z^{-6}$.

  \item Find all cube roots of
  \[
  -8i
  \]
  and give the answers in both polar form and standard form.
\end{enumerate}

\end{document}
